\section*{序論}
11月2日朝から3日朝にかけての生成AI関連X投稿では、モデルを単に推論に使う段階から、別のモデルや開発環境を操作する段階への広がりが目立った。特に、ニューラルネットワークの重みを直接読み書きするという研究報告が関心を集めた。一方、Mistral Large 4の技術レポート公開を通じ、オープンウェイトモデルの実装と性能を巡る議論も続いている。

開発ツールの側では、コーディングエージェントのセッション記憶、仕様駆動、プラグイン生成、共有ストレージなど、継続的な作業を支える仕組みが複数紹介された。AI学習のオプトアウト信号を巡る議論は、技術的な合意と実際の提供時期の差を浮かび上がらせた。

以下はGrokが収集した10件を原文の順序で記事化したものである。未確認の発言や性能値は投稿・資料上の主張として扱い、独立したWeb検証は実施していない。
